\documentclass[%
%  reprint,
 %preprint,
% onecolumn,
    twocolumn,
superscriptaddress,
%groupedaddress,
%unsortedaddress,
%runinaddress,
%frontmatterverbose, 
preprintnumbers,
nofootinbib,
%nobibnotes,
%bibnotes,
 amsmath,amssymb,
 aps,
 prd,
 floatfix,
]{revtex4-2}

% WIDTHS per \the\columnwidth
% Column: 246pt / 72.27pt/in = 3.4039 in
% Widetext: 510pt / 72.27pt/in = 7.05687 in

\usepackage{graphicx}% Include figure files
\usepackage{duckuments} % DEBUG: removeme
\usepackage{dcolumn}% Align table columns on decimal point
\usepackage{bm}% bold math
\usepackage{braket}
\usepackage{mathtools} % ensuremath etc
\usepackage{enumitem} % for bullet dedentation
\usepackage{slashed}

\usepackage[pdftex]{xcolor}
\usepackage{hyperref}% add hypertext capabilities
\hypersetup{
    colorlinks=true,     % false: boxed links; true: colored links
    linkcolor=blue,      % color of internal links
    citecolor=blue,      % color of links to bibliography
    filecolor=blue,      % color of file links
    urlcolor=blue        % color of external links
}

\usepackage{cleveref}
\crefname{equation}{Eq.}{Eqs.}
\crefname{figure}{Fig.}{Figs.}
\crefname{table}{Tab.}{Tabs.}
\crefname{section}{Sec.}{Secs.}
\crefname{appendix}{App.}{Apps.}
\Crefname{table}{Table}{Tables}
\Crefname{figure}{Figure}{Figures}

\newcommand{\tx}[0]{\text{TX}}
\newcommand{\qcd}[0]{\text{QCD}}
\newcommand{\txqcd}[0]{\text{TXQCD}}

%%% NOTATION MACROS
\newcommand{\aset}[1]{\{#1\}}
\newcommand{\RRvec}[0]{\xi}
\newcommand{\RRmat}[0]{E}
\newcommand{\Pm}[1]{\mathcal{P}^{(#1)}}
\newcommand{\HEq}[0]{\rightarrow}
\newcommand{\Span}{\mathrm{span}}
\newcommand{\herm}{\mathfrak{h}}

%%% EDITING MACROS
\definecolor{red}{rgb}{1.0, 0, 0}
\newcommand{\draftnote}[1]{{\color{red} {#1}}}
\newcommand{\todo}[1]{{\color{red} \textbf{TODO:} {#1}}}
\newcommand{\dan}[1]{{\color{purple} {#1}}}
\newcommand{\ryan}[1]{{\color{cyan} {#1}}}
\newcommand{\mike}[1]{{\color{magenta} {#1}}}
\newcommand{\dimitra}[1]{{\color{olive} {#1}}}


%%% CONVENIENCE MACROS

\begin{document}

%\preprint{FERMILAB-PUB-25-0131-T,MIT-CTP/5849}

% \title{Coincidences in Krylov Space\\(or, Rayleigh-Ritz is all you need)}
\title{Notes on tensor extended QCD}

% \author{Ryan Abbott}
%  \affiliation{Center for Theoretical Physics, Massachusetts Institute of Technology, Cambridge, MA 02139, USA}
%  \affiliation{The NSF AI Institute for Artificial Intelligence and Fundamental Interactions}
% \author{Daniel C. Hackett}
%  \affiliation{Fermi National Accelerator Laboratory, Batavia, IL 60510, U.S.A.}
% \author{George T.\ Fleming}
%  \affiliation{Fermi National Accelerator Laboratory, Batavia, IL 60510, U.S.A.}
% \author{Dimitra A. Pefkou}
%  \affiliation{Department of Physics, University of California, Berkeley, CA 94720, USA}
%  \affiliation{Nuclear Science Division, Lawrence Berkeley National Laboratory, Berkeley, CA 94720, USA}
 \author{Michael L. Wagman}
  \affiliation{Fermi National Accelerator Laboratory, Batavia, IL 60510, U.S.A.}



%\begin{abstract}
%     We introduce a new twist on Kaplan's extended QCD proposal~\cite{Kaplan:2013dca}.
%\end{abstract}


\maketitle


\section{Introduction}
\label{sec:intro}

Kaplan proposed eXtended QCD (XQCD) in Ref.~\cite{Kaplan:2013dca} in particular as a way improve the signal-to-noise problem arising for nucleon correlators based on successful applications to NJL models~\cite{Grabowska:2012ik,Nicholson:2012xt}.
A scalar auxiliary field picks up a VEV that shifts the quark propagator poles, therefore effectively adding a ``constituent quark mass.''
At the same time, pions are kept light through the presence of disconnected diagrams induced by mixing with a flavored pseudoscalar auxiliary field.
The additional attraction arising from the presence of this auxiliary scalar is compensated by repulsion arising from auxiliary vector and axial-vector fields.
The balance between the different auxiliary fields, and therefore exact equivalence with QCD, is guaranteed by adding auxiliary fields associated with quark bilinears whose products sum to zero by a Fierz identity.


Unfortunately, quark determinants in XQCD are complex---the SNR problem is traded for a sign problem. These notes explore a variation on XQCD that involves antisymmetric-tensor auxiliary fields and preserves reality of the quark determinant.

\section{Tensor XQCD}

Tensor XQCD (TXQCD) is defined by adding a new term to the QCD Lagrangian in Euclidean spacetime,
\begin{equation}\label{eq:LTX}
\begin{split}
   \mathcal{L}_{\tx} &= \frac{\lambda^2}{2} \left[ \left(\sigma_{ab} + \frac{ \overline{q}^i_a q^i_b }{\lambda^2} \right)\left(\sigma_{ba} + \frac{ \overline{q}^j_b q^j_a }{\lambda^2} \right) \right. \\
   &\hspace{20pt} + \left(\pi_{ab} + \frac{ \overline{q}^i_a \gamma_5 q^i_b }{\lambda^2} \right)\left(\pi_{ba} + \frac{ \overline{q}^j_b \gamma_5 q^j_a }{ \lambda^2} \right) \\
   &\hspace{20pt} + \left(s^{ij} + \frac{ \overline{q}^i_a q^j_a }{\sqrt{2} \lambda^2} \right)\left(s^{ji} + \frac{ \overline{q}^j_b q^i_b }{\sqrt{2} \lambda^2} \right) \\
    &\hspace{20pt} + \left(p^{ij} + \frac{ \overline{q}^i_a \gamma_5 q^j_a }{\sqrt{2} \lambda^2} \right)\left(p^{ji} + \frac{ \overline{q}^j_b \gamma_5 q^i_b }{\sqrt{2} \lambda^2} \right) \\
    &\hspace{20pt} \left. + \left(t_{\mu\nu}^{ij} + \frac{ \overline{q}^i_a \sigma_{\mu\nu} q^j_a }{2 \lambda^2} \right)\left(t_{\mu\nu}^{ji} + \frac{ \overline{q}^j_b \sigma_{\mu\nu} q^i_b }{2 \lambda^2} \right) \right],
   \end{split}
\end{equation}
where the auxiliary fields $\sigma_{ab}(x)$, $\pi_{ab}(x)$, $s^{ij}(x)$, $p^{ij}(x)$ and $t^{ij}_{\mu\nu}(x) = -t^{ij}_{\nu\mu}(x)$ are Hermitian matrices in either flavor $a,b,\ldots$ or color $i,j,\ldots$ indices, $q_a^i$ and $\overline{q}_a^i$  are quark and antiquark fields with spinor indices implicit, and $\sigma_{\mu\nu} \equiv (i/2)[\gamma_{\mu},\gamma_{\nu}]$  is Hermitian, $\sigma_{\mu\nu} = \sigma_{\mu\nu}^\dagger$.

The parameter $\lambda$ and all auxiliary fields have mass dimension one.

The Fierz identity
\begin{equation}
\begin{split}
    0 &= (\overline{q}^i_a q^j_b)(\overline{q}^k_c  q^l_d) + (\overline{q}^i_a  \gamma_5 q^j_b)(\overline{q}^k_c  \gamma_5 q^l_d)  \\
    &\hspace{10pt} + \frac{1}{2} (\overline{q}^i_a q^l_d)(\overline{q}^k_c q^j_b) + \frac{1}{2}  (\overline{q}^i_a \gamma_5 q^l_d)(\overline{q}^k_c \gamma_5 q^j_b) \\
    &\hspace{10pt} + \frac{1}{4} (\overline{q}^i_a \sigma_{\mu\nu} q^l_d)(\overline{q}^k_c \sigma_{\mu\nu} q^j_b) ,
    \end{split}
\end{equation}
which follows from
\begin{equation}
\begin{split}
    &\delta_{
    \alpha\beta} \delta_{\gamma\delta} + [\gamma_5]_{\alpha\beta}[\gamma_5]_{\gamma\delta} \\
    &\hspace{10pt} = \frac{1}{2} \delta_{\alpha\delta} \delta_{\gamma\beta} + \frac{1}{2} [\gamma_5]_{\alpha\delta}[\gamma_5]_{\gamma\beta} + \frac{1}{4} [\sigma_{\mu\nu}]_{\alpha\delta}[\sigma_{\mu\nu}]_{\gamma\beta},
    \end{split}
\end{equation}
shows that the quartic quark interactions in $\mathcal{L}_{\tx}$ vanish. This provides the equivalent expression
\begin{equation}
\begin{split}
    \mathcal{L}_{\tx} &= \frac{\lambda^2}{2}\text{Tr}
    [\bm{\sigma}^2 + \bm{\pi}^2 + s^2 + p^2 + t_{\mu\nu} t_{\mu\nu}] \\
    & + \overline{q}^i_a\left[ \sigma_{ab} + \pi_{ab}\gamma_5 + \frac{s^{ij}}{\sqrt{2}}  + \frac{p^{ij}}{\sqrt{2}}  \gamma_5 + \frac{t_{\mu\nu}^{ij}}{2}  \sigma_{\mu\nu}    \right]q^j_b,
    \end{split}
\end{equation}
where boldface symbols denote $N_f \times N_f$ flavor matrices and $s$, $p$, and $t_{\mu\nu}$ are implicitly $N_c \times N_c$ color matrices.

Path integrals of $S_{\tx} \equiv \sum_x \mathcal{L}_{\tx}$ are Gaussian and can be performed analytically e.g. by shifting the real (imaginary) part of each auxiliary field component by ($-i$ times) the symmetric (antisymmetric) part of the corresponding quark bilinear matrix in \cref{eq:LTX},
\begin{equation}\label{eq:aux}
\begin{split}
    \int \mathcal{D}\sigma\mathcal{D}\pi\mathcal{D}s\mathcal{D}p\mathcal{D}t_{\mu\nu} \ e^{-S_{\tx}} = \mathcal{N}, \\
    \end{split}
\end{equation}
where $\mathcal{N}$ is a normalization constant that depends on $\lambda$ but is independent of the quark field.
By first performing the auxiliary field integrals, observables in TXQCD can be seen to be identical to those in QCD by construction
\begin{equation}
\begin{split}
\left< \mathcal{O} \right>_{\txqcd} &\equiv \frac{1}{Z \mathcal{N}} \int \mathcal{D}(...)\mathcal{D}\overline{q} \mathcal{D}q\ e^{-S_{\qcd} - S_{\tx}} \mathcal{O} \\
&= \frac{1}{Z} \int \mathcal{D}U \mathcal{D}\overline{q} \mathcal{D}q\ e^{-S_{\qcd}} \mathcal{O} \\
&= \left< \mathcal{O} \right>_{\qcd},
\end{split}
\end{equation}
where $\mathcal{D}(...)$ is $\mathcal{D}U$ times the auxiliary field path integral measure in \cref{eq:aux} and $Z$ is the usual QCD partition function.
Here, it is assumed that $\mathcal{O}$ is constructed from products of QCD quark and gluon operators and does not contain auxiliary fields explicitly, although more general expectation values can be physically interpreted~\cite{Kaplan:2013dca}.


\subsection{Quark determinant reality}

Instead integrating out the quarks first gives
\begin{equation}\label{eq:TXQCD}
\begin{split}
\left< \mathcal{O}_{\txqcd} \right> 
&= \frac{1}{Z} \int \mathcal{D}(...) \ e^{-S_g} e^{-\frac{\lambda^2}{2} \sum_x \text{Tr}[\bm{\sigma}^2+ \bm{\pi}^2 + s^2 + p^2 + t_{\mu\nu}t_{\mu\nu} ]} \\
&\hspace{20pt} \times \det(\slashed{D}_{\txqcd} + m),
\end{split}
\end{equation}
where $S_g$ is the QCD gauge action and
\begin{equation}
\begin{split}
     \bm{\slashed{D}}_{\txqcd} &= \slashed{D}_{\qcd} + \bm{\sigma} + \bm{\pi} \gamma_5   \\
    &\hspace{10pt} + \frac{s}{\sqrt{2}} + \frac{p}{\sqrt{2}} \gamma_5 + \frac{1}{2} t_{\mu\nu}\sigma_{\mu\nu}.
    \end{split}
\end{equation}

Suppose the QCD Dirac operator $\slashed{D}_{\qcd}$ satisfies $\gamma_5 \slashed{D}_{\qcd} \gamma_5 = \slashed{D}_{\qcd}^\dagger$.
Then
\begin{equation}
\begin{split}
    &\gamma_5 \slashed{D}_{\txqcd} \gamma_5 \\
    &= \slashed{D}_{\qcd}^\dagger + \bm{\sigma} + \bm{\pi} \gamma_5  + \frac{s}{\sqrt{2}} + \frac{p}{\sqrt{2}} \gamma_5 + \frac{1}{2} t_{\mu\nu}\sigma_{\mu\nu} \\
    &= \slashed{D}_{\qcd}^\dagger + \left( \bm{\sigma} + \bm{\pi} \gamma_5  + \frac{s}{\sqrt{2}} + \frac{p}{\sqrt{2}} \gamma_5 + \frac{1}{2} t_{\mu\nu}\sigma_{\mu\nu} \right)^\dagger \\
    &= \slashed{D}_{\txqcd}^\dagger.
    \end{split}
\end{equation}
It follows that $\det(\slashed{D}_{\txqcd} + m)$ is real.

It can be further shown that the eigenvalues of $\slashed{D}_{\txqcd}$ are either real or come in complex conjugate pairs.
To see this, assume that $\mu$ is an eigenvalue of $\slashed{D}_{\txqcd}$ and therefore a root of the characteristic equation, $\det(\mu  - \slashed{D}_{\txqcd}) = 0$.
This means
\begin{equation}
    \begin{split}
        & \det(\mu -  \slashed{D}_{\txqcd}^\dagger ) \\
        &=  \det(\mu  - \gamma_5 \slashed{D}_{\txqcd} \gamma_5 ) \\
        &= \det(\gamma_5) \det(\mu  - \slashed{D}_{\txqcd} ) \det( \gamma_5 ) \\
        &= \det(\mu - \slashed{D}_{\txqcd}) \\
        &= 0.
    \end{split}
\end{equation}
It follows that $\det \slashed{D}_{\txqcd}$ is positive iff $\slashed{D}_{\txqcd}$ has an even number of exactly real, negative eigenvalues.

In practice, positivity can be verified as with strange quark determinants in QCD: perform simulations with $|\det(\slashed{D}_{\txqcd})|$, compute the lowest eigenvalues of the Hermitian operator
\begin{equation}
    \gamma_5 \slashed{D}_{\txqcd} = ( \gamma_5 \slashed{D}_{\txqcd})^\dagger, 
\end{equation}
for instance with the Lanczos algorithm,
and verify that the smallest eigenvalue is positive.
If so, this is sufficient to guarantee that $\det(\gamma_5 \slashed{D}_{\txqcd}) = \det( \slashed{D}_{\txqcd}) \geq 0$ for a given field configuration.


\subsection{VEVs and mean-field picture}

As in the original XQCD setup~\cite{Kaplan:2013dca}, the scalar auxiliary fields get VEVs
\begin{equation}
    \lambda^2 \left< \sigma_{ab} \right> = -\delta_{ab} \frac{1}{N_f} \left< \overline{q}^i_c q^i_c \right> \equiv \delta_{ab} \Sigma,
\end{equation}
and
\begin{equation}
    \lambda^2 \left< s^{ij} \right> = -\delta^{ij} \frac{1}{\sqrt{2} N_c} \left< \overline{q}^k_a q^k_a \right> = \delta^{ij} \frac{N_f}{\sqrt{2} N_c} \Sigma.
\end{equation}
Within a mean-field approximation for the gluon and auxiliary fields, this shifts the pole in the fermion propagator as $m \rightarrow m + M$, where the generated ``constituent quark mass'' $M$ is
\begin{equation}\label{eq:M_Sigma}
\begin{split}
    M =  \left< \frac{\text{Tr}[\bm{\sigma}]}{N_f}   + \frac{\text{Tr}[s]}{\sqrt{2} N_c}  \right>  = \frac{\Sigma}{\lambda^2}\left( 1 + \frac{N_f}{2N_c} \right).
    \end{split}
\end{equation}
Quark propagators, whose magnitude in QCD decays with a correlation length set by the inverse of $m_\pi /2$, will have their decay rate set by $m_\pi/2 + M$. 

If $\lambda$ is chosen so that $3M \sim M_N - (3/2) m_\pi$, then the magnitudes of baryon correlation functions will decay at a rate set by $M_N$ rather then the usual $(3/2)m_\pi$.
Within this mean-field picture, the quark propagator products entering nucleon correlators in TXQCD are $\sim e^{-M_N t}$, exponentially closer to average baryon correlator values than the $\sim e^{-(3/2)m_\pi t}$ arising in QCD.
Phase fluctuations are therefore not required to contribute a significant amount of the nucleon mass in TXQCD correlators, unlike in QCD where they contribute $M_N - (3/2)m_\pi$~\cite{Wagman:2016bam}.
Since the variance of a phase is $O(1)$, these phase fluctuations lead to the $M_N - (3/2)m_\pi$ exponential SNR degradation of nucleon correlators.
Removing the need for phase fluctuations to contribute most of the nucleon mass should correspondingly improve the SNR problem.

If $M \sim \Lambda_{\rm QCD}$ is achieved, then based off the empirical positivity of strange-quark determinants one expects one may expect positivity of the TXQCD determinant to be achieved as well.

What about pion correlators? Pions $\overline{q}\tau^A \gamma_5 q$ couple to the $\bm{\pi}$ auxiliary field, and therefore pion correlators include disconnected contributions that can produce light pion poles as in Ref.~\cite{Kaplan:2013dca}. 
This allows the existence of pseudo-Golstone pion poles to be reconciled with the presence of an $O(\Lambda_{\rm QCD})$ constituent quark mass.

\section{TXQCD$_2$}

In two dimensions, the analog of the chiral-even Fierz relation above is 
\begin{equation}
\begin{split}
    0 &= (\overline{q}^i_a q^j_b)(\overline{q}^k_c  q^l_d) + (\overline{q}^i_a  \sigma_3 q^j_b)(\overline{q}^k_c  \sigma_3 q^l_d)  \\
    &\hspace{10pt} + (\overline{q}^i_a q^l_d)(\overline{q}^k_c q^j_b) + (\overline{q}^i_a \sigma_3 q^l_d)(\overline{q}^k_c \sigma_3 q^j_b) ,
    \end{split}
\end{equation}
which follows from
\begin{equation}
\begin{split}
    &\delta_{
    \alpha\beta} \delta_{\gamma\delta} + [\sigma_3]_{\alpha\beta}[\sigma_3]_{\gamma\delta} = \delta_{\alpha\delta} \delta_{\gamma\beta} +  [\sigma_3]_{\alpha\delta}[\sigma_3]_{\gamma\beta},
    \end{split}
\end{equation}
allows us to construct an analogous auxiliary field action for QCD$_2$,
\begin{equation}\label{eq:LTX2}
\begin{split}
   \mathcal{L}_{\tx_2} &= \frac{\lambda^2}{2} \left[ \left(\sigma_{ab} + \frac{ \overline{q}^i_a q^i_b }{\lambda} \right)\left(\sigma_{ba} + \frac{ \overline{q}^j_b q^j_a }{\lambda} \right) \right. \\
   &\hspace{20pt} + \left(\pi_{ab} + \frac{ \overline{q}^i_a \sigma_3 q^i_b }{\lambda^2} \right)\left(\pi_{ba} + \frac{ \overline{q}^j_b \sigma_3 q^j_a }{ \lambda} \right) \\
   &\hspace{20pt} + \left(s^{ij} + \frac{ \overline{q}^i_a q^j_a }{\lambda} \right)\left(s^{ji} + \frac{ \overline{q}^j_b q^i_b }{\lambda} \right) \\
    &\hspace{20pt} \left. + \left(p^{ij} + \frac{ \overline{q}^i_a \sigma_3 q^j_a }{\lambda} \right)\left(p^{ji} + \frac{ \overline{q}^j_b \sigma_3 q^i_b }{\lambda} \right) \right],
   \end{split}
\end{equation}
which has the equivalent form
\begin{equation}
\begin{split}
    \mathcal{L}_{\tx_2}  &= \frac{\lambda^2}{2}\text{Tr}
    [\bm{\sigma}^2 + \bm{\pi}^2 + s^2 + p^2] \\
    & \hspace{10pt} + \lambda \overline{q}^i_a\left[ \sigma_{ab} + \pi_{ab}\sigma_3 + s^{ij} + p^{ij}  \sigma_3   \right]q^j_b.
   \end{split}
\end{equation}
As in the four-dimensional case, auxiliary field integrals of $e^{-S_{\tx_2}}$, where $S_{\tx_2} \equiv \sum_x \mathcal{L}_{\tx_2}$, are Gaussian and can be performed analytically to show that TXQCD$_2$ observables defined by the action $S_{\qcd_2} + S_{\txqcd_2}$ are identical to QCD$_2$ observables after auxiliary field integration.
In the two-dimensional case, $\lambda$ still has mass dimension one but the auxiliary fields have mass dimension zero, as appropriate for canonically normalized scalars with $D=2$.


Instead integrating out the quarks first gives
\begin{equation}\label{eq:TXQCD2}
\begin{split}
\left< \mathcal{O}_{\txqcd_2} \right> 
&= \frac{1}{Z} \int \mathcal{D}(...) \ e^{-S_g} e^{-\frac{\lambda^2}{2} \sum_x \text{Tr}[\bm{\sigma}^2+ \bm{\pi}^2 + s^2 + p^2 ]} \\
&\hspace{20pt} \times \det(\slashed{D}_{\txqcd_2} + m),
\end{split}
\end{equation}
where $S_g$ is the QCD$_2$ gauge action and
\begin{equation}
\begin{split}
  \bm{\slashed{D}}_{\txqcd_2} &= \slashed{D}_{\qcd_2} + \lambda\left( \bm{\sigma} + \bm{\pi} \sigma_3  + s + p \sigma_3 \right).
    \end{split}
\end{equation}


Suppose the QCD$_2$ Dirac operator $\slashed{D}_{\qcd_2}$ satisfies $\sigma_3 \slashed{D}_{\qcd_2} \sigma_3 = \slashed{D}_{\qcd_2}^\dagger$.
Then
\begin{equation}
\begin{split}
    &\sigma_3 \slashed{D}_{\txqcd} \sigma_3 \\
    &= \slashed{D}_{\qcd_2}^\dagger + \lambda \left(\bm{\sigma} + \bm{\pi} \sigma_3  + s + p \sigma_3\right) \\
    &= \slashed{D}_{\qcd_2}^\dagger + \lambda \left( \bm{\sigma} + \bm{\pi} \sigma_3  + s + p \sigma_3 \right)^\dagger \\
    &= \slashed{D}_{\txqcd_2}^\dagger.
    \end{split}
\end{equation}
It follows that $\det(\slashed{D}_{\txqcd_2} + m)$ is real.

As in the four-dimensional case, the scalar fields get VEVs
\begin{equation}
    \lambda \left< \sigma_{ab} \right> = -\delta_{ab} \frac{1}{N_f} \left< \overline{q}^i_c q^i_c \right> \equiv \delta_{ab} \Sigma,
\end{equation}
and
\begin{equation}
    \lambda \left< s^{ij} \right> = -\delta^{ij} \frac{1}{N_c} \left< \overline{q}^k_a q^k_a \right> = \delta^{ij} \frac{N_f}{N_c} \Sigma.
\end{equation}
The corresponding constituent quark mass is
\begin{equation}\label{eq:M_Sigma_2}
\begin{split}
    M =  \lambda \left< \frac{\text{Tr}[\bm{\sigma}]}{N_f}   + \frac{\text{Tr}[s]}{N_c}  \right>  = \Sigma \left( 1 + \frac{N_f}{N_c} \right).
    \end{split}
\end{equation}

The is a critical difference between the role of $\lambda$ in TXQCD and TXQCD$_2$.
Upon rescaling all auxiliary fields as $\bm{\sigma} \rightarrow \lambda \bm{\sigma}$, $\bm{\pi} \rightarrow \lambda \bm{\pi}$, $s \rightarrow \lambda s$, and $p\rightarrow \lambda p$, the TXQCD$_2$ action becomes independent of $\lambda$.
This means that not only expectation values but also quark propagators and configuration-by-configuration values of correlation functions will be independent of $\lambda$.
The parameter $\lambda$ cannot affect the variances of observables or any other practical aspect of Monte Carlo calculations of TXQCD$_2$.
Numerical simulations may set $\lambda = 1$ for convenience with no need to study different values of $\lambda$ besides as a consistency check on the implementation.
Unfortunately, this means that there is no possibility to optimize SNR by tuning $\lambda$; the result for any $\lambda > 0$ is universal.
Note however that $\lambda = 0$ is a singular point where the auxiliary field path integrals diverge and therefore it is consistent for configuration-by-configuration values of e.g. quark propagators to differ between QCD$_2$ (where $M$ vanishes) and TXQCD$_2$ with $\lambda > 0$.

\section{One-site integral equations}

\subsection{Two dimensions}

The equivalence of QCD$_2$ and XQCD$_2$ implies the determinant relation
\begin{equation}
  \begin{split}
  \det(A) &= \int \mathcal{D}\sigma \mathcal{D}\pi \mathcal{D}s \mathcal{D}p \ e^{-\frac{\lambda^2}{2} \sum_x \text{Tr}[\bm{\sigma}^2+ \bm{\pi}^2 + s^2 + p^2 ]} \\
  &\hspace{10pt} \times \det(A  + \bm{\sigma} + \bm{\pi} \sigma_3  + s + p \sigma_3 )
  \end{split}
\end{equation}
is true for any tensor field $A_{ab\alpha\beta}^{ij}(x,y)$ with $\alpha,\beta \in \{1,2\}$, $i,j\in \{1,\ldots,N_c\}$, and $a,b\in \{1,\ldots,N_f\}$.
%2N+4*(N-1)
%=6N-4
For a single lattice site, this becomes a $(6N_f+6N_c - 8)$-dimensional integral equation,
\begin{equation}
  \begin{split}
    &\det(A_{ab\alpha\beta}^{ij}) = \int_{-\infty}^{\infty} \prod_a d\sigma_{aa}d\pi_{aa} \prod_{b>a} d\sigma_{ab}^R d\sigma_{ab}^I d\pi_{ab}^R d\pi_{ab}^I \\ 
    &\hspace{40pt}\times \prod_i ds^{ii}dp^{ii} \prod_{j>i} ds^{ij}_R ds^{ij}_I dp^{ij}_R dp^{ij}_I \\
    &\hspace{40pt} \times  e^{-\frac{\lambda^2}{2} \sum_{a,b} \left( |\sigma_{ab}|^2 + |\pi_{ab}|^2 \right)} e^{-\frac{\lambda^2}{2} \sum_{i,j} \left( |s^{ij}|^2 + |p^{ij}|^2 \right) } \\
    &\hspace{40pt} \times \det\left[A_{ab\alpha\beta}^{ij}  + \left(\sigma_{ab}\delta_{\alpha\beta} + \pi_{ab} [\sigma_3]_{\alpha\beta}\right)\delta^{ij} \right. \\
    &\hspace{70pt} \left. + \left(s^{ij}\delta_{\alpha\beta} + p^{ij} [\sigma_3]_{\alpha\beta}\right)\delta_{ab} \vphantom{A_{ab\alpha\beta}^{ij}}\right],
  \end{split}
\end{equation}
with all indices and integral measures written out fully explicitly and $R$ ($I$) subscripts/superscripts denote real (imaginary) parts.
We have verified this integral relation by performing the Gaussian integrals on the right-hand-side analytically for a few choices of $A$, $N_f$, and $N_c$.

For $N_c = 1$ or $N_f = 1$, valid integral relations are obtained when restricting the auxiliary fields to be purely real and symmetric. For $N_c, N_f \geq 2$ it is necessary to include imaginary parts that are antisymmetric in the relevant matrix indices in order to obtain valid relations.
Does this reflect the fact that $(\overline{u}^i u^j)$ and $(\overline{u}^i \sigma_3 u^j)$ are symmetric under $i\leftrightarrow j$ but $(\overline{q}^i_a q^j_b)$ is not?
In the latter case, antisymmetric matrix components are required in order to obtain Gaussian auxiliary field integerals with means that describe all components of $(\overline{q}^i_a q^j_b)$; however, the Fierz identities appear to involve only $(\overline{q}^i_a q^i_b)$ and $(\overline{q}^i_a q^j_a)$ so it's not obvious to me where this restriction arises from.

The result for $A_{ab\alpha\beta}^{ij} = \delta_{ab}\delta_{\alpha\beta}\delta^{ij}$ gives the normalization constant
\begin{equation}
    \mathcal{N} = 2^{N_f + N_c} \left( \frac{\pi}{\lambda^2} \right)^{3N_f + 3N_c - 4 },
\end{equation}
which for the full $2D$ theory is simply raised to the $V$ power, where $V$ is the number of lattice sites.

The result for $A_{ab\alpha\beta}^{ij} = a^2 m \delta_{ab}\delta_{\alpha\beta}\delta^{ij}$ divided by the free quark partition function $(a^2 m)^{2N_f N_c}$ gives the result for $N_f$ flavors of static Wilson-like quarks with masses $m a \gg 1$,
\begin{equation}
  \begin{split}
    \Sigma &= \frac{\lambda}{N_f}\left< \text{Tr}\bm{\sigma} \right> = \frac{\lambda}{N_f} \left< \text{Tr}s \right> \\
    &= \frac{2 N_c}{m a^2} \left[ 1 + O\left(\frac{1}{m a}\right) \right].
  \end{split}
\end{equation}
This is consistent with a direct calculation using the QCD$_2$ quark propagator,
\begin{equation}
  \begin{split}
    -\frac{1}{N_f}\left< \overline{q}^i_a q^i_a \right>_{\qcd_2} &= \frac{1}{N_f} \left< \text{Tr}\left[ \slashed{D}_{\qcd_2}^{-1} \right] \right> \\
    &= \frac{2 N_c}{m a^2} \left[ 1 + O\left(\frac{1}{m a}\right) \right],
  \end{split}
\end{equation}
since $[\slashed{D}_{\qcd_2}]_{\alpha\beta ab}^{ij} = a^2 m \delta_{\alpha\beta}\delta_{ab}\delta^{ij}$ plus kinetic and other terms suppressed by $1/(ma)$.
It is interesting to note that although a direct calculation using the \txqcd$_2$ propagator is also guaranteed to give the same answer, such a calculation is significantly more complicated to perform than the calculation of $\left< \text{Tr}\bm{\sigma} \right>$
\begin{equation}
  \begin{split}
    &-\frac{1}{N_f}\left< \overline{q}^i_a q^i_a \right>_{\txqcd_2} = \frac{1}{N_f} \left< \text{Tr}\left[ \slashed{D}_{\txqcd_2}^{-1} \right] \right> \\
    &= \frac{1}{\mathcal{N} (a^2 m)^{2N_f N_c}} \int_{-\infty}^{\infty} \prod_a d\sigma_{aa}d\pi_{aa} \prod_{b>a} d\sigma_{ab}^R d\sigma_{ab}^I d\pi_{ab}^R d\pi_{ab}^I \\ 
    &\hspace{40pt}\times \prod_i ds^{ii}dp^{ii} \prod_{j>i} ds^{ij}_R ds^{ij}_I dp^{ij}_R dp^{ij}_I \\
    &\hspace{40pt} \times  e^{-\frac{\lambda^2}{2} \sum_{a,b} \left( |\sigma_{ab}|^2 + |\pi_{ab}|^2 \right)} e^{-\frac{\lambda^2}{2} \sum_{i,j} \left( |s^{ij}|^2 + |p^{ij}|^2 \right) } \\
    &\hspace{40pt} \times \det\left[a^2 m \delta^{ij}\delta_{\alpha\beta}\delta_{ab}  + \left(\sigma_{ab}\delta_{\alpha\beta} + \pi_{ab} [\sigma_3]_{\alpha\beta}\right)\delta^{ij} \right. \\
    &\hspace{70pt} \left. + \left(s^{ij}\delta_{\alpha\beta} + p^{ij} [\sigma_3]_{\alpha\beta}\right)\delta_{ab} \vphantom{A_{ab\alpha\beta}^{ij}}\right] \\
    &\hspace{40pt} \times \text{Tr}\left[a^2 m \delta^{ij}\delta_{\alpha\beta}\delta_{ab}  + \left(\sigma_{ab}\delta_{\alpha\beta} + \pi_{ab} [\sigma_3]_{\alpha\beta}\right)\delta^{ij} \right. \\
    &\hspace{70pt} \left. + \left(s^{ij}\delta_{\alpha\beta} + p^{ij} [\sigma_3]_{\alpha\beta}\right)\delta_{ab} \vphantom{A_{ab\alpha\beta}^{ij}}\right],
  \end{split}
\end{equation}
because the inverse Dirac operator is a rational, rather than polynomial, function of the auxiliary fields.

The constituent quark mass in this limit is
\begin{equation}
  \begin{split}
    M &= \frac{\lambda}{N_f}\left< \text{Tr}\bm{\sigma} \right> + \frac{\lambda}{N_c} \left< \text{Tr} s \right> \\
    &= \frac{2(N_c + N_f)}{m a^2} \left[ 1 + O\left(\frac{1}{ma}\right) \right].
  \end{split}
\end{equation}
Importantly, the constituent quark mass is equal to a positive constant times the bare quark mass $m$ in this limit, meaning that the inclusion of TXQCD$_2$ auxiliary fields increases the decay rate of quark propagators in free field theory.

As discussed above, specifically in $D=2$ all TXQCD$_2$ observables are independent of $\lambda$.
This means that $M$ takes an action-dependent value that cannot be tuned by changing $\lambda$.
This is guaranteed to be the case even outside the large $ma$ limit by \cref{eq:M_Sigma_2}, $M = \Sigma (1 + N_f/N_c)$.
It can also be seen from the fact that $\lambda$ can be removed from the theory entirely be rescaling all auxiliary fields by $\lambda$, e.g. $s \rightarrow \lambda s$.


\subsection{Four dimensions}

In the four-dimensional case, equivalence of QCD and TXQCD implies
\begin{equation}
  \begin{split}
  \det(A) &= \int \mathcal{D}\sigma \mathcal{D}\pi \mathcal{D}s \mathcal{D}p \mathcal{D}t_{\mu\nu} \ e^{-\frac{\lambda^2}{2} \sum_x \text{Tr}[\bm{\sigma}^2+ \bm{\pi}^2 + s^2 + p^2 + t_{\mu\nu}t_{\mu\nu}]} \\
  &\hspace{10pt} \times \det\left(A  + \bm{\sigma} + \bm{\pi} \gamma_5  + \frac{s}{\sqrt{2}} + \frac{p \gamma_5 }{\sqrt{2}} + \frac{t^{ij}_{\mu\nu} \sigma_{\mu\nu} }{2} \right),
  \end{split}
\end{equation}
is valid for any $A_{ab\alpha\beta}^{ij}(x,y)$ with $\alpha,\beta \in \{1,\ldots,4\}$, $i,j\in~\{1,\ldots,N_c\}$, and $a,b\in \{1,\ldots,N_f\}$.
%8N+16*(N-1)
%=24N-16
For a single lattice site, this becomes a $(6N_f+24N_c - 20)$-dimensional integral equation,
\begin{equation}
  \begin{split}
    &\det(A_{ab\alpha\beta}^{ij}) = \int_{-\infty}^{\infty} \prod_a d\sigma_{aa}d\pi_{aa} \prod_{b>a} d\sigma_{ab}^R d\sigma_{ab}^I d\pi_{ab}^R d\pi_{ab}^I \\ 
    &\hspace{40pt}\times \prod_i ds^{ii}dp^{ii} \prod_{j>i} ds^{ij}_R ds^{ij}_I dp^{ij}_R dp^{ij}_I \\
     &\hspace{40pt}\times \prod_{\mu} \prod_{\nu > \mu} \prod_i dt_{\mu\nu}^{ii} \prod_{j>i} d[t_{\mu\nu}^{ij}]_R d[t_{\mu\nu}^{ij}]_I  \\
    &\hspace{40pt} \times  e^{-\frac{\lambda^2}{2} \sum_{a,b} \left( |\sigma_{ab}|^2 + |\pi_{ab}|^2 \right)} e^{-\frac{\lambda^2}{2} \sum_{i,j} \left( |s^{ij}|^2 + |p^{ij}|^2 \right) } \\
    &\hspace{40pt} \times e^{-\lambda^2 \sum_{i,j} \sum_{\mu < \nu} |t_{\mu\nu}^{ij}|^2  } \\
    &\hspace{40pt} \times \det\left[A_{ab\alpha\beta}^{ij}  + \left(\sigma_{ab}\delta_{\alpha\beta} + \pi_{ab} [\gamma_5]_{\alpha\beta}\right)\delta^{ij} \right. \\
    &\hspace{70pt} \left. + \frac{1}{\sqrt{2}} \left(s^{ij}\delta_{\alpha\beta} + p^{ij} [\gamma_5]_{\alpha\beta} \right)\delta_{ab} \right. \\
    &\hspace{70pt} \left. + \frac{1}{2} \sum_{\mu,\nu} t^{ij}_{\mu\nu}[\sigma_{\mu\nu}]_{\alpha\beta} \delta_{ab}  \vphantom{A_{ab\alpha\beta}^{ij}}\right].
  \end{split}
\end{equation}
We have again verified this integral relation analytically for a few choices of $A$, $N_f$, and $N_c$.
%As in the $D=2$ case, complex Hermitian auxiliary field matrices are required to obtain valid identities when $N_c, N_f \geq 2$.

The result for $A_{ab\alpha\beta}^{ij} = \delta_{ab}\delta_{\alpha\beta}\delta^{ij}$ gives the normalization constant
\begin{equation}
    \mathcal{N} = 2^{N_f - 5N_c + 6} \left( \frac{\pi}{\lambda^2} \right)^{3N_f + 12N_c - 10 },
\end{equation}
which for the full $4D$ theory is simply raised to the $V$ power, where $V$ is the number of lattice sites.

The result for $A_{ab\alpha\beta}^{ij} = a^4 m \delta_{ab}\delta_{\alpha\beta}\delta^{ij}$  divided by the free quark partition function $(a^4 m)^{4N_f N_c}$ can be used to compute the chiral condensate for $D=4$ LQCD with $N_f$ static Wilson-like quarks with masses $m a \gg 1$,
\begin{equation}
  \begin{split}
    \Sigma &= \frac{\lambda^2}{N_f}\left< \text{Tr}\bm{\sigma} \right> = \frac{\sqrt{2} \lambda^2}{N_f} \left< \text{Tr}s \right> \\
    &= \frac{4 N_c}{m a^4} \left[ 1 + O\left(\frac{1}{m a}\right) \right].
  \end{split}
\end{equation}
This is consistent with a direct calculation using the QCD quark propagator,
\begin{equation}
  \begin{split}
    -\frac{1}{N_f}\left< \overline{q}^i_a q^i_a \right>_{\qcd} &= \frac{1}{N_f} \left< \text{Tr}\left[ \slashed{D}_{\qcd}^{-1} \right] \right> \\
    &= \frac{4 N_c}{m a^4} \left[ 1 + O\left(\frac{1}{m a}\right) \right],
  \end{split}
\end{equation}
since $[\slashed{D}_{\qcd}]_{\alpha\beta ab}^{ij} = a^4 m \delta_{\alpha\beta}\delta_{ab}\delta^{ij}$ plus kinetic and other terms suppressed by $1/(ma)$.
As in the two-dimensional case, direct calculation using the TXQCD quark propagator trace must give the same answer but is much more cumbersome than TXQCD auxiliary field expectation value above.

The constituent quark mass is
\begin{equation}
  \begin{split}
    M &= \frac{1}{N_f}\left< \text{Tr}\bm{\sigma} \right> + \frac{1}{\sqrt{2} N_c} \left< \text{Tr} s \right> \\
    &= \frac{4 N_c}{\lambda^2 m a^4}\left( 1 + \frac{N_f}{2N_c} \right) \left[ 1 + O\left(\frac{1}{m a}\right) \right].
  \end{split}
\end{equation}
This is consistent with \cref{eq:M_Sigma}, $\lambda^2 M = \Sigma (1 + N_f/(2N_c))$.
As in the two-dimensional case, the sign of $M$ is identical to the sign of $m$ and the decay rate of quark propagators is increased in free field theory.
Unlike the two-dimensional case, $M$ is proportional to $\lambda^{-2}$ and can be tuned to any desired value by changing $\lambda$.
There is a one-parameter family of TXQCD theories in $D=4$ that give identical expectation values for QCD fields but can have different variances of e.g. products of quark propagators viewed as functions of the gluon field.  

\bibliography{Lanczos_bib}

\end{document}
